\section{时频局部化：Fourier 丢掉了``何时''}
\label{sec:08-Numerical-Analysis/08-Fourier-and-Spectral-Methods/05}

一台机组的振动记录，两分钟，采样率 \(20\) kHz。运维报告说其中某个时刻转速跳了一次，要你从数据里把这个时刻找出来。

把整段数据做一次 FFT，谱上确实出现了两个主峰，对应两个转速。但谱图上没有任何东西告诉你哪个峰在前、哪个在后，甚至没告诉你它们是不是一直同时存在。构造两段人工信号就能看清这件事有多彻底。在 \([0,T]\) 上取

\[
f_1(t)=
\begin{cases}
\sin(2\pi\nu_1 t), & 0\le t<T/2,\\[2pt]
\sin(2\pi\nu_2 t), & T/2\le t\le T,
\end{cases}
\qquad
f_2(t)=\frac1{\sqrt2}\bigl(\sin(2\pi\nu_1 t)+\sin(2\pi\nu_2 t)\bigr).
\]

\(f_1\) 先低后高，每个频率各占一半时间；\(f_2\) 两个频率从头到尾同时存在。作为工程对象，一个是转速切换，一个是双频共振，处理方式完全不同。可两者的幅度谱在 \(\nu_1\) 与 \(\nu_2\) 处各有一个主峰，量级相当；由 Parseval 恒等式，总能量也相同。把 \(T\) 取得更长不会改善，只会让两个峰更尖。

Fourier 变换本身没丢信息，它是可逆的，时间顺序完整地留在相位里。麻烦在于相位没法直接读——工程上画的是 \(\abs{X_k}\)，而顺序恰好躺在被扔掉的那一半里。更根本的原因在基函数：\(e^{2\pi i\nu t}\) 在整条时间轴上永恒振荡，它自身不携带任何位置信息。把信号投到这组基上，等于问一首曲子用了哪些音符：答案精确，可先后顺序被平均掉了。旋律和把所有音符同时砸下的一个和弦，在这份清单上难以分辨。

\subsection{切开再变换，时间就回来了}

既然整段平均是病因，最直接的修复是先切开。用一个以 \(\tau\) 为中心的窗口 \(w\) 乘上去，只对截出的那一小段做变换：

\[
G_wf(\tau,\omega)\eqdef\int_{-\infty}^{\infty}f(t)\,\overline{w(t-\tau)}\,e^{-i\omega t}\,dt .
\]

结果有两个变量，\(\tau\) 回答何时，\(\omega\) 回答什么频率。这就是短时 Fourier 变换（short-time Fourier transform, STFT），把 \(\abs{G_wf(\tau,\omega)}^2\) 画在时频平面上就是谱图（spectrogram）。对上面的 \(f_1\)，谱图前半段亮在低频、后半段亮在高频，转速切换的时刻一眼可见；对 \(f_2\)，两条亮带从头贯穿到尾。

实现上没有新东西。以固定步长滑动窗口，对每帧 \(N\) 个样本做一次\hyperref[sec:08-Numerical-Analysis/08-Fourier-and-Spectral-Methods/03]{``DFT 是变换，FFT 是算法''}里的 DFT，把各帧的幅度谱沿时间轴堆起来。窗口不必是矩形，\hyperref[sec:08-Numerical-Analysis/08-Fourier-and-Spectral-Methods/02]{``采样、混叠与可恢复性''}里讨论频谱泄漏时出现的 Hann、Blackman 窗在每一帧内照旧适用，主瓣与旁瓣的那笔账逐帧重算一遍。语音处理里几乎所有声学前端都停在这一步：\(25\) ms 的窗、\(10\) ms 的帧移，做完 STFT 再压到梅尔刻度上，就是喂给模型的那张梅尔谱图。

于是 STFT 可以概括成 Fourier 变换加一道局部化预处理。它的全部力量来自那扇窗，全部麻烦也是。

\subsection{窗宽一选定，两个分辨率就绑死了}

窗口一引入就得定宽度，而两个方向各有各的代价。窗窄，时间定位准，能把紧挨着的两个事件分开；可窗内装不下几个周期，频率就算不准——可分辨的频率间隔大致是窗长的倒数，\(25\) ms 的窗给出 \(40\) Hz 上下的频率分辨率，低音区的相邻半音就此糊在一起。窗宽，频率分辨细，可窗内发生的一切都被压到同一个 \(\tau\) 附近，时间被抹平，一次持续 \(5\) ms 的瞬态在 \(100\) ms 的窗里只是一小片背景。

把这件事画出来最清楚。每个窗口在时频平面上占一块宽约 \(\Delta t\)、高约 \(\Delta\omega\) 的分辨率格子，STFT 就是拿一批形状全同的格子去铺满这个平面。

\begin{center}
\begin{tikzpicture}[scale=1.0,font=\small,>=stealth]
  \begin{scope}
    \draw[->,black!55] (-0.1,0) -- (4.6,0) node[right,font=\footnotesize] {\(t\)};
    \draw[->,black!55] (0,-0.1) -- (0,4.5) node[above,font=\footnotesize] {\(\omega\)};
    \foreach \i in {0,1,2,3,4,5,6,7}{
      \foreach \j in {0,1,2,3}{
        \draw[blue!55!black,fill=blue!8] (\i*0.5,\j*1.0) rectangle (\i*0.5+0.5,\j*1.0+1.0);
      }
    }
    \draw[<->,black!70] (0.05,4.28) -- (0.45,4.28);
    \node[above,font=\footnotesize,black!70] at (0.25,4.3) {\(\Delta t\)};
    \draw[<->,black!70] (4.22,2.05) -- (4.22,2.95);
    \node[right,font=\footnotesize,black!70] at (4.24,2.5) {\(\Delta\omega\)};
    \node[font=\footnotesize,black!70] at (2.0,-0.55) {窄窗：时间看得清};
  \end{scope}
  \begin{scope}[shift={(6.6,0)}]
    \draw[->,black!55] (-0.1,0) -- (4.6,0) node[right,font=\footnotesize] {\(t\)};
    \draw[->,black!55] (0,-0.1) -- (0,4.5) node[above,font=\footnotesize] {\(\omega\)};
    \foreach \i in {0,1,2,3}{
      \foreach \j in {0,1,2,3,4,5,6,7}{
        \draw[orange!70!black,fill=orange!8] (\i*1.0,\j*0.5) rectangle (\i*1.0+1.0,\j*0.5+0.5);
      }
    }
    \draw[<->,black!70] (0.05,4.28) -- (0.95,4.28);
    \node[above,font=\footnotesize,black!70] at (0.5,4.3) {\(\Delta t\)};
    \draw[<->,black!70] (4.22,2.05) -- (4.22,2.45);
    \node[right,font=\footnotesize,black!70] at (4.24,2.25) {\(\Delta\omega\)};
    \node[font=\footnotesize,black!70] at (2.0,-0.55) {宽窗：频率看得清};
  \end{scope}
\end{tikzpicture}

\emph{两种铺法的格子面积一样，只是长宽比不同。换窗宽就是在同一块面积里重新分配时间与频率，而面积本身压不下去。}
\end{center}

图里那句``面积压不下去''是有定量含义的。用

\[
\Delta t^2=\frac{\int t^2\abs{w(t)}^2\,dt}{\int\abs{w(t)}^2\,dt},
\qquad
\Delta\omega^2=\frac{\int \omega^2\abs{\hat w(\omega)}^2\,d\omega}{\int\abs{\hat w(\omega)}^2\,d\omega}
\]

度量窗口在两个域上的展宽（假设都已中心化），则

\[
\Delta t\,\Delta\omega\ge\frac12 .
\]

\subsection{这个下界是变换对的结构性质}

证明的骨架只有三步。先对 \(\int\abs{w}^2\) 做分部积分，把 \(1\) 写成 \(t\) 的导数，得到

\[
\begin{aligned}
\norm{w}^2&=-\int t\,\frac{d}{dt}\abs{w(t)}^2\,dt=-2\operatorname{Re}\int t\,w'(t)\,\overline{w(t)}\,dt,\\
\norm{w}^2&\le 2\norm{t\,w}\cdot\norm{w'},
\end{aligned}
\]

第二步是 Cauchy--Schwarz。第三步用 Plancherel 把 \(\norm{w'}\) 换成 \(\norm{\omega\hat w}\)（差一个与约定有关的常数），整理即得 \(\Delta t\,\Delta\omega\ge1/2\)。

条件对账：分部积分要求边界项 \(t\abs{w(t)}^2\) 在无穷远处趋于零，这对任何有限能量且衰减合理的窗都成立，但对不衰减的窗（例如把整段信号原样取来的``无窗''）就不成立，那时 \(\Delta t\) 本来就是无穷；Cauchy--Schwarz 取等号当且仅当 \(w'\) 与 \(t w\) 成正比，解这个微分方程给出 \(w(t)=Ce^{-\alpha t^2}\)，也就是说 Gaussian 窗是唯一取到下界的窗，其余的窗都只能做得更差。常数 \(1/2\) 依赖频率约定，用 \(\nu\) 代替 \(\omega\) 会多出 \(2\pi\) 因子，这与\hyperref[sec:08-Numerical-Analysis/08-Fourier-and-Spectral-Methods/01]{``从 Fourier 级数到连续变换''}里那条约定警告是同一件事，但正下界的存在与约定无关。

所以窄窗与宽窗之间的取舍不是实现不够好，也不是没找到聪明的窗形。它写在 Fourier 变换对的结构里：时域压紧，频域必定摊开。工程上唯一的自由度是选在哪一档，不是能不能不选。

\subsection{全同格子对不上信号的真实需求}

铺全同格子还有一个隐含后果：所有频率被同一个分辨率审视。而信号在不同频段上想要的东西并不一样。

高频成分变化快，一个周期很短，值得用窄窗把时间钉准；频率轴上少分辨几赫兹，相对误差也不大——\(8\) kHz 的成分测偏 \(40\) Hz 是千分之五。低频成分变化慢，一个周期横跨很长时间，\(25\) ms 的窗装不下一个 \(30\) Hz 的完整周期，给它宽窗才能测出频率，损失一点时间定位无所谓，低频事件本来就没有那么陡的时间结构。同样是偏 \(40\) Hz，对 \(60\) Hz 的成分就是灾难。

固定窗宽被迫在这两种需求之间做一次全局妥协，而且妥协的结果对整个频段一视同仁。实践里常见的补救是同时算几幅不同窗宽的谱图对照着看——这个做法本身就承认了单一分辨率不够用。时间序列建模里滑窗长度的选择、多尺度特征的拼接，遇到的是同一个困难的另一种说法：窗口一固定，能看见的变化速度就固定了。

出路是让窗宽随频率走：分析高频时自动收缩，分析低频时自动放宽，时频平面上的铺法从全同矩形变成分层的：高频那几层的格子窄而高，低频那几层宽而矮。固定窗做不到这一点，因为它只有一个宽度参数可调。下一节把``缩放''从外部参数提升为基函数自身的组织方式，那张铺法图会被重画一次。
